\section{因子分析法}
\subsection{授業内容}
\subsubsection{科目の中でのこのコマの位置づけ}
目に見えないものを測定するという意味で，テスト理論は心理学の測定と関係が深い。前回は社会心理学における態度の測定を前提に議論したが，測定に関してはテスト理論で一般的に議論することができる。
真のスコアと誤差とに分解すること，誤差の基本的な仮定を確認した上で，古典的テスト理論を項目と被験者の特性に分割することで因子分析モデルに展開されるところを見る。また，多因子モデルに拡張した上で，その数理的展開から，信頼性と妥当性に言及できることを確認する。数式の展開は代数の基本的な特徴を確認すれば問題なくフォローできるだろう。
\subsubsection{コマ主題細目}
\begin{description}
    \item[古典的テスト理論]
    \item[因子分析モデル]
    \item[因子分析の第2定理]
    \item[因子分析の第1定理]
\end{description}
\subsubsection{キーワード}
\begin{itemize}
\item 古典的テスト理論
\item 因子分析法
\item 因子分析の定理
\end{itemize}
\subsection{授業情報}
\paragraph{コマの展開方法}
講義
\subsubsection{予習・復習課題}
\paragraph{予習}
\paragraph{復習}
